\begin{figure*}[t]
\centering
\begin{tikzpicture}[node distance=0.48cm, every node/.style={font=\scriptsize}]
  \node[dmblock, text width=2.25cm] (src) {Prompt CSVs\\baseline + strategies};
  \node[dmblockgreen, text width=2.25cm, right=of src] (merge) {experiment\\Prompt table\\\code{prompt_uid}};
  \node[dmblockorange, text width=2.05cm, right=of merge] (api) {DF gateway\\12 models};
  \node[dmblock, text width=2.15cm, right=of api] (raw) {raw response\\long table};
  \node[dmblockgreen, text width=2.15cm, right=of raw] (feat) {behavior\\feature table};
  \node[dmblockorange, text width=2.05cm, right=of feat] (mine) {analysis\\tables};
  \draw[dmarrow] (src) -- (merge);
  \draw[dmarrow] (merge) -- (api);
  \draw[dmarrow] (api) -- (raw);
  \draw[dmarrow] (raw) -- (feat);
  \draw[dmarrow] (feat) -- (mine);
  \node[dmnote, below=0.42cm of api, text width=6.0cm] {断点续跑键：$(\mcode{provider},\mcode{model},\mcode{repeat},\mcode{prompt_uid})$};
  \node[dmnote, below=0.42cm of feat, text width=5.0cm] {失败样本保留 \code{status} 与 \code{error_message}，不静默丢弃。};
\end{tikzpicture}
\caption{多模型行为数据采集管线。上游 Prompt 文件首先合并为带有全局唯一键的实验主表，随后经 DF/OpenAI-compatible 网关调用多个模型，最终形成同时保留原始回复、结构化特征和失败记录的长表数据。}
\label{fig:collection-pipeline}
\end{figure*}
